%!TEX root = ../../main.tex

\chapter{Erstellung des synthetischen Datensatzes}
\label{ch:datensatz}

Das in Kapitel~\ref{ch:konzeption} entworfene Untersuchungsdesign setzt einen Datensatz voraus, dessen Eigenschaften exakt bekannt sind: Nur wenn die enthaltenen Datenqualitätsprobleme gezielt eingebracht wurden und für jede Aufgabe eine eindeutige Soll-Lösung existiert, lassen sich die in Abschnitt~\ref{sec:evaluationskriterien} definierten Kriterien überhaupt objektiv messen. Dieses Kapitel dokumentiert daher die Erstellung des synthetischen Datensatzes. Er wird einmalig vorab erzeugt und über sämtliche Läufe unverändert verwendet. Seine Erzeugung erfolgt in zwei Schritten, indem zunächst eine vollständig saubere, schemakonforme Fassung entsteht und diese anschließend kontrolliert verfälscht wird. Es schließt unmittelbar an die in Abschnitt~\ref{sec:etl-aufgaben} definierten neun \ac{ETL}-Aufgaben an und zeigt, wie genau die dort geforderten Probleme in kontrolliertem Umfang eingebracht werden.

Zunächst werden die aus dem Untersuchungsdesign abgeleiteten Anforderungen begründet (Abschnitt~\ref{sec:datensatz-anforderungen}), anschließend das Datenmodell und die saubere Zielstruktur beschrieben (Abschnitt~\ref{sec:datenmodell}). Darauf folgt die detaillierte Zuordnung der gezielt eingebauten Qualitätsprobleme zu den einzelnen Aufgaben (Abschnitt~\ref{sec:eingebaute-probleme}), bevor das zweistufige Generierungsverfahren (Abschnitt~\ref{sec:generierung}) und die abschließende Validierung des Datensatzes (Abschnitt~\ref{sec:datensatz-validierung}) dargelegt werden. Der so erzeugte Datensatz bildet die Eingabe für die technische Umsetzung in Kapitel~\ref{ch:implementierung} und die Durchführung der Experimente in Kapitel~\ref{ch:durchfuehrung}.

\section{Anforderungen an den Datensatz}
\label{sec:datensatz-anforderungen}

Aus den Forschungsfragen und dem in Kapitel~\ref{ch:konzeption} festgelegten Design ergeben sich vier Anforderungen:

\noindent\textbf{Repräsentativität:} Der Datensatz muss die in Abschnitt~\ref{subsec:dq-probleme} beschriebenen, praxistypischen Datenqualitätsprobleme abbilden, und zwar gerade jene, die den neun in Abschnitt~\ref{sec:etl-aufgaben} definierten Aufgaben entsprechen, sodass jede Aufgabe an einem tatsächlich vorhandenen Mangel ansetzt \autocite{ridzuanReviewDataQuality2024,coteDataCleaningMachine2024}.

\noindent\textbf{Kontrollierbarkeit:} Sollzustand sowie Art und Umfang der eingebrachten Fehler müssen exakt bekannt sein, da die Verarbeitungsqualität gegen eine bekannte Soll-Lösung gemessen wird. Bei realen Beständen wäre diese Referenz nur mit erheblichem Annotationsaufwand und verbleibender Unsicherheit herstellbar. Die Wahl regelbasiert erzeugter Testdaten ist eine Setzung dieses Versuchsdesigns. Sie ermöglicht die Kontrolle der eingebrachten Muster, belegt aber keine Repräsentativität für reale Datenbestände.

\noindent\textbf{Reproduzierbarkeit:} Die Generierung muss deterministisch sein, damit die Datengrundlage keine Variabilität beisteuert, die der querschnittlich erhobenen Reproduzierbarkeit der Modellergebnisse zugerechnet würde.

\noindent\textbf{Angemessener Umfang:} Der Datensatz muss groß genug sein, um die drei Schwierigkeitsstufen jeder Aufgabe auszuprägen, und klein genug, um die in Abschnitt~\ref{subsec:effizienz} erhobenen Aufwandskennzahlen nicht durch reine Datenmenge zu dominieren.

\section{Datenmodell und -struktur}
\label{sec:datenmodell}

Als Anwendungsszenario dient ein bewusst einfach gehaltenes E-Commerce-Datenmodell, das aus den drei Tabellen Kunden, Produkte und Bestellungen besteht. Diese Wahl ist kein Selbstzweck, sondern folgt aus den Aufgaben: Das Szenario ist allgemein verständlich, deckt sämtliche in Abschnitt~\ref{sec:etl-aufgaben} geforderten Problemtypen ab und stellt über die Beziehung der Bestellungen zu Kunden und Produkten zugleich die Schlüsselstruktur bereit, die für die anspruchsvolle Transformationsaufgabe aus Verknüpfung, Bereinigung und Aggregation benötigt wird \autocite{khanOverviewETLTechniques2024,walhaDataIntegrationTraditional2024}. Die Tabelle der Bestellungen verweist über Fremdschlüssel auf Kunden und Produkte und verbindet die drei Tabellen so zu einem kleinen relationalen Modell.

Grundlage des Modells ist die Definition eines sauberen Zielschemas für jede Tabelle, das den fehlerfreien Sollzustand festlegt. Für die Kundentabelle umfasst dieses Schema eine eindeutige Kundennummer, den vollständigen Namen, die E-Mail-Adresse, einen zweistelligen Länder-Code nach ISO-3166-1-alpha-2 sowie das Registrierungsdatum. Die Produkttabelle besteht aus einer Produktnummer, der Bezeichnung, einer Kategorie, dem Preis als nichtnegativem Dezimalwert und einem Wahrheitswert zur Verfügbarkeit. Die Bestelltabelle schließlich verknüpft Bestell-, Kunden- und Produktnummer mit der bestellten Menge, dem Einzelpreis zum Bestellzeitpunkt und dem Bestelldatum. Der Gesamtbetrag ist dort bewusst \emph{nicht} enthalten, da seine Berechnung Gegenstand der mittleren Transformationsaufgabe ist. Er tritt erst in deren Soll-Lösung als abgeleitete Spalte hinzu.

Dieses saubere Schema erfüllt eine doppelte Funktion: Es beschreibt einerseits das von den Modellen zu erreichende Zielformat und dient andererseits, in seiner konkret erzeugten Ausprägung, als bekannte Referenz für die Bewertung. Die typisierte, maschinenlesbare Festlegung des Schemas stellt dabei sicher, dass die als Referenz erzeugten Daten strikt schemakonform sind und damit als verlässliche Ground Truth taugen \autocite{ridzuanReviewDataQuality2024,rangineniReviewEnhancingData2023}.

Der Umfang des Datensatzes ist so gewählt, dass er die genannte Abwägung zwischen Vielfalt und Verarbeitungsaufwand einlöst. Die Kundentabelle umfasst in ihrer sauberen Fassung fünfhundert Datensätze, die Produkttabelle einhundert und die Bestelltabelle zweitausend Datensätze. Diese Größenordnung ist eine Setzung des Versuchsdesigns: groß genug, um Qualitätsprobleme merklich einzustreuen und die Aggregation über Länder und Kategorien belastbar zu füllen, klein genug für die wiederholte Ausführung über alle Kombinationen. Ob sie trägt, belegen erst die Aufwandskennzahlen in Abschnitt~\ref{sec:aufwand}.

\section{Eingebaute Datenqualitätsprobleme}
\label{sec:eingebaute-probleme}

Kern der Datensatzkonstruktion ist das gezielte Einbringen von Qualitätsproblemen, das dieser Abschnitt entlang der drei in Abschnitt~\ref{sec:etl-aufgaben} definierten Aufgabenkategorien beschreibt. Jedes Problem ist dabei nicht beliebig gewählt, sondern korrespondiert mit genau einer Aufgabe und einer der in Abschnitt~\ref{subsec:dq-dimensionen} eingeführten Qualitätsdimensionen, sodass sich die spätere Bewertung lückenlos auf einen bewusst gesetzten Mangel zurückführen lässt \autocite{ridzuanReviewDataQuality2024,coteDataCleaningMachine2024}. Auf diese Weise ist die in Abschnitt~\ref{sec:etl-aufgaben} eingeführte Matrix aus neun Aufgaben und drei Schwierigkeitsstufen unmittelbar in der Struktur des Datensatzes verankert.

Die Probleme der \emph{Datenbereinigung} sind ausschließlich in der Kundentabelle verortet und in aufsteigender Schwierigkeit gestaffelt: zusätzliche Leerzeichen in den Namensfeldern und entfernte Länderangaben für die einfache Stufe, das in vier Darstellungen abgelegte Registrierungsdatum für die mittlere und die Verfälschung der Länderangaben durch wechselnde Schreibweisen, Sprachen, Abkürzungen und gezielte Tippfehler für die anspruchsvolle \autocite{coteDataCleaningMachine2024,rangineniReviewEnhancingData2023}. Gerade die letztgenannte Verfälschung bildet die in Abschnitt~\ref{subsec:aufgabe-bereinigung} hervorgehobene semantische Achse ab, da ihre Auflösung ein inhaltliches Verständnis der Werte und nicht nur einen syntaktischen Abgleich erfordert \autocite{hassaniRoleChatGPTData2023,khanCognitiveETLUsing2025}.

Die Probleme der \emph{Datentransformation} betreffen die Produkt- und die Bestelltabelle und knüpfen an die Dimension der Konsistenz an. In der Produkttabelle werden die numerischen Preisangaben in uneinheitliche textuelle Schreibweisen überführt und die Verfügbarkeit von einem Wahrheitswert in wechselnde textuelle Entsprechungen umgeschrieben, womit die einfache Aufgabe der Typkonvertierung ihre Grundlage erhält. Die mittlere Stufe benötigt keine eigene Verfälschung, sondern nutzt die vorhandene Struktur der Bestelltabelle, aus der sich Gesamtbetrag und Datumsbestandteile ableiten lassen. Die anspruchsvolle Stufe schließlich verlangt die Verknüpfung aller drei Tabellen über ihre Schlüsselbeziehungen, eine begleitende Bereinigung der zuvor beschriebenen Mängel und eine abschließende Aggregation über Land und Produktkategorie. Sie bündelt damit bewusst mehrere eingebaute Probleme in einem mehrstufigen Verarbeitungsschritt, wie er für reale \ac{ETL}-Strecken charakteristisch ist \autocite{khanOverviewETLTechniques2024,walhaDataIntegrationTraditional2024}.

Die Probleme der \emph{Duplikaterkennung} zielen auf die Eindeutigkeit der Kundentabelle. In die saubere Ausgangsbasis von 500 Kunden werden für die einfache und mittlere Stufe 20 vollständig identische Datensätze eingefügt, die sowohl als exakte Duplikate als auch über die mehrfach vergebene E-Mail-Adresse erkennbar sind. Für die anspruchsvolle Stufe treten weitere 20 Datensätze als unscharfe Duplikate (Fuzzy-Duplikate) mit abweichender Schreibweise von Name und E-Mail hinzu. Die verfälschte Kundentabelle (\texttt{customers\_raw.csv}) umfasst somit insgesamt 540 Zeilen. Die Soll-Lösungen der drei Duplikataufgaben unterscheiden sich entsprechend im Umfang. Die einfache und die mittlere Stufe verlangen jeweils nur die Entfernung der 20 exakten beziehungsweise schlüsselgleichen Datensätze, ihre Soll-Lösung umfasst folglich $540 - 20 = 520$ Datensätze. Die anspruchsvolle Stufe zielt dagegen auf die tatsächlich eindeutigen Personen und verlangt damit die Auflösung beider Duplikatklassen, sodass ihre Soll-Lösung $540 - 20 - 20 = 500$ Datensätze umfasst, die 500 eindeutigen Kunden des Sollzustands. Dieselben 500 Datensätze werden auch in der durchgängigen Bereinigung und in der verketteten Pipeline erreicht, dort jedoch über zwei nacheinander ausgeführte Schritte \autocite{coteDataCleaningMachine2024,hassaniRoleChatGPTData2023}. Der Anteil der eingebrachten Probleme ist so bemessen, dass er die Ergebniskennzahlen merklich beeinflusst, ohne den Charakter des Datensatzes als überwiegend sauberen Bestand zu verfälschen.

\begin{table}[htbp]
    \centering
    \small
    \caption[Beispiele verfälschter Werte im Datensatz]{Beispiele tatsächlich erzeugter Rohwerte und ihrer Sollwerte, gegliedert nach der zugehörigen Aufgabe. Sichtbare Leerzeichen sind mit \textvisiblespace{} gekennzeichnet.}
    \label{tab:dirty-werte}
    \begin{tabularx}{\textwidth}{@{}l>{\raggedright\arraybackslash}X>{\raggedright\arraybackslash}X@{}}
        \toprule
        \textbf{Aufgabe} & \textbf{Rohwert} & \textbf{Sollwert} \\
        \midrule
        Leerzeichen &\texttt{\textvisiblespace\textvisiblespace Sigfried Römer-Butte\textvisiblespace\textvisiblespace} & \texttt{Sigfried Römer-Butte} \\
        fehlender Wert & \texttt{} (leer) & \texttt{UNKNOWN} \\
        \addlinespace
        Datumsformat & \texttt{2024-12-11} (ISO) & \texttt{2024-12-11} \\
        Datumsformat & \texttt{02.01.2023} (deutsch) & \texttt{2023-01-02} \\
        Datumsformat & \texttt{August 09 2025} (US) & \texttt{2025-08-09} \\
        Datumsformat & \texttt{1628121600} (Unix) & \texttt{2021-08-05} \\
        \addlinespace
        Länderangabe & \texttt{Deutschland} / \texttt{Germany} & \texttt{DE} \\
        Länderangabe & \texttt{Deutshcland} (Tippfehler) & \texttt{DE} \\
        Länderangabe & \texttt{SUI} / \texttt{Schweiz} & \texttt{CH} \\
        Länderangabe & \texttt{Holland} / \texttt{Niederlande} & \texttt{NL} \\
        Länderangabe & \texttt{England} / \texttt{UK} & \texttt{GB} \\
        \addlinespace
        Preisangabe & \texttt{494.35 EUR} & \texttt{494.35} \\
        Preisangabe & \texttt{\texteuro 6.53} & \texttt{6.53} \\
        Preisangabe & \texttt{155,17} & \texttt{155.17} \\
        Verfügbarkeit & \texttt{yes} / \texttt{true} / \texttt{1} & \texttt{True} \\
        Verfügbarkeit & \texttt{no} / \texttt{false} / \texttt{0} & \texttt{False} \\
        \bottomrule
    \end{tabularx}
\end{table}

Wie sich das konkret in den Daten niederschlägt, zeigt Tabelle~\ref{tab:dirty-werte} anhand tatsächlicher Zeilen des erzeugten Datensatzes. Die Gegenüberstellung von Roh- und Sollwert macht zugleich sichtbar, weshalb die Aufgaben unterschiedlich schwer sind: Während sich Leerzeichen und Datumsformate durch eine endliche Menge syntaktischer \newline Regeln vollständig auflösen lassen, ist bei den Länderangaben für jeden einzelnen Wert zu entscheiden, welches Land gemeint ist. Fehlende Werte werden nicht gelöscht, sondern einheitlich als \texttt{UNKNOWN} gekennzeichnet, sodass die Zeilenzahl erhalten bleibt und die Vollständigkeit im Sinne von Abschnitt~\ref{subsec:vollstaendigkeit} messbar bleibt.

Da der Kern eines Duplikats im Verhältnis zweier Datensätze zueinander liegt, stellt Tabelle~\ref{tab:dirty-duplikate} sie zeilenweise gegenüber. Bei den exakten Duplikaten ist der eingefügte Datensatz in allen Feldern einschließlich der Kundennummer identisch. Bei den unscharfen weichen Name und E-Mail-Adresse durch genau jene Muster ab, die auch in realen Beständen entstehen, veränderte Groß- und Kleinschreibung, ein zusätzliches Leerzeichen, ein Punkt im lokalen Teil der Adresse. Die inhaltliche Identität ist bei manueller Sichtung erkennbar, ist über einen exakten Vergleich jedoch nicht mehr zugänglich, weshalb diese Aufgabe die in Abschnitt~\ref{subsec:aufgabe-duplikate} beschriebene Ähnlichkeitsbetrachtung erzwingt.

\begin{table}[htbp]
    \centering
    \small
    \caption[Beispiele eingefügter Duplikate]{Eingefügte Duplikate im Vergleich zu ihrem Ausgangsdatensatz. Bei den unscharfen Duplikaten sind die abweichenden Stellen die einzige Änderung gegenüber dem Original.}
    \label{tab:dirty-duplikate}
    \begin{tabularx}{\textwidth}{@{}llX X@{}}
        \toprule
        \textbf{Art} & \textbf{Nr.} & \textbf{Name} & \textbf{E-Mail} \\
        \midrule
        exakt (Original)  & 21 & \texttt{Janina Pieper-Junken} & \texttt{patriziakarge@example.com} \\
        exakt (Duplikat)  & 21 & \texttt{Janina Pieper-Junken} & \texttt{patriziakarge@example.com} \\
        \addlinespace
        unscharf (Original) & 49 & \texttt{Herr Alban Kusch} & \texttt{janeroskoth@example.net} \\
        unscharf (Duplikat) & 503 & \texttt{Herr\textvisiblespace\textvisiblespace Alban\textvisiblespace\textvisiblespace Kusch} & \texttt{\textvisiblespace janeroskoth@example.net\textvisiblespace} \\
        \addlinespace
        unscharf (Original) & 494 & \texttt{Dr. Käthe Mülichen} & \texttt{doris08@example.net} \\
        unscharf (Duplikat) & 501 & \texttt{DR. KÄTHE MÜLICHEN} & \texttt{DORIS08@example.net} \\
        \addlinespace
        unscharf (Original) & 391 & \texttt{Nadine Butte} & \texttt{shertrampf@example.org} \\
        unscharf (Duplikat) & 508 & \texttt{nadine butte} & \texttt{shert.rampf@example.org} \\
        \bottomrule
    \end{tabularx}
\end{table}

Aus den 500 sauberen Kundendatensätzen entstehen durch das Anfügen von 20 exakten und 20 unscharfen Duplikaten insgesamt 540 Rohdatensätze, was die in Abschnitt~\ref{sec:datenmodell} genannte Zielgröße von 500 eindeutigen Kunden als Sollzustand der Duplikataufgaben festlegt. Die Häufigkeitsverteilung der Länderangaben bestätigt zugleich die beabsichtigte Streuung: Für die zehn im Datensatz vertretenen Länder liegen je Land drei bis sechs verschiedene Schreibweisen vor, von Standardcodes wie \texttt{DE} über ausgeschriebene deutsche und englische Namen bis zu untypischen Kürzeln wie \texttt{SUI} für die Schweiz oder \texttt{GER} für Deutschland, und 38 Angaben fehlen vollständig.

\section{Generierungsverfahren und verwendete Werkzeuge}
\label{sec:generierung}

Die Erzeugung des Datensatzes folgt einem zweistufigen Prinzip, das die in Abschnitt~\ref{sec:datensatz-anforderungen} geforderte Kontrollierbarkeit unmittelbar umsetzt. Im ersten Schritt werden für jede Tabelle vollständig saubere, strikt schemakonforme Daten erzeugt, die den fehlerfreien Sollzustand verkörpern. Im zweiten Schritt werden aus dieser sauberen Fassung durch eine kontrollierte Verfälschung die tatsächlich zu verarbeitenden Rohdaten abgeleitet, während die saubere Fassung selbst als Referenz erhalten bleibt. Dieses Vorgehen kehrt den üblichen Ablauf einer \ac{ETL}-Strecke bewusst um: Nicht die Bereinigung wird nachträglich auf unbekannte Rohdaten angewandt, sondern die Verunreinigung wird gezielt auf bekannte saubere Daten aufgebracht, sodass die exakte Umkehrung jeder Verfälschung und damit die eindeutige Soll-Lösung jeder Aufgabe von vornherein feststeht.

Für die Erzeugung realistisch wirkender Basiswerte wie Namen und E-Mail-Adressen wird die Bibliothek \texttt{Faker} eingesetzt, deren deutschsprachige Lokalisierung dem Szenario angemessene Werte liefert. Die kontrollierte Verfälschung, also das Einstreuen von Leerzeichen, fehlenden Werten, uneinheitlichen Formaten, verfälschten Länderangaben sowie exakten und unscharfen Duplikaten, erfolgt anschließend durch eigens implementierte Routinen, die jede der in Abschnitt~\ref{sec:eingebaute-probleme} beschriebenen Problemklassen erzeugen. Sowohl die Basisdaten als auch die Verfälschung werden über einen zentralen Zufallszahlengenerator gesteuert, der mit einem festen Startwert initialisiert wird. Dadurch ist die Generierung vollständig deterministisch: Derselbe Startwert führt stets zu einem identischen Datensatz, während unterschiedliche Startwerte voneinander unabhängige, aber strukturell gleichartige Datensätze erzeugen. Die erzeugten Rohdaten werden in einem gängigen tabellarischen Austauschformat abgelegt und dienen zusammen mit der separat gespeicherten sauberen Referenz als Eingabe für die in Kapitel~\ref{ch:implementierung} beschriebene Verarbeitungsumgebung.

\section{Validierung des Datensatzes}
\label{sec:datensatz-validierung}

Der Datensatz wird vor seiner Verwendung programmatisch geprüft. Ein Validierungsskript erzeugt ihn dazu aus dem hinterlegten Startwert neu und führt vier Prüfungen aus, von denen jede fehlschlagen und den Datensatz damit verwerfen kann. Die erste validiert die saubere Referenz zeilenweise gegen die typisierte Schemadefinition aus Abschnitt~\ref{sec:datenmodell}: Datentyp, Wertebereich und Format jedes Feldes müssen der Festlegung entsprechen \autocite{ridzuanReviewDataQuality2024,rangineniReviewEnhancingData2023}. Nur unter dieser Voraussetzung ist der spätere Abgleich der Modellergebnisse mit der Referenz aussagekräftig.

Die zweite Prüfung weist für jede der neun Aufgaben aus Abschnitt~\ref{sec:etl-aufgaben} den Mangel nach, an dem sie ansetzt, und zählt ihn aus: Randleerzeichen und fehlende Werte, die Zahl der auftretenden Datums- und Preisformate, die Zahl der vom Zielschema abweichenden Länderangaben sowie mehrfach vorkommende und unscharf abgewandelte Datensätze. Die dritte stellt sicher, dass zu jeder Aufgabe eine nicht leere Soll-Lösung vorliegt, was die in Abschnitt~\ref{sec:evaluationskriterien} definierten Kriterien überhaupt erst berechenbar macht \autocite{coteDataCleaningMachine2024}. Die vierte erzeugt den Datensatz ein zweites Mal und vergleicht beide Fassungen bytegenau, was die Reproduzierbarkeit der Generierung bestätigt. Für den verwendeten Datensatz bestehen sämtliche Prüfungen.

Mit dieser Validierung steht ein Datensatz zur Verfügung, dessen Sollzustand, Fehlerbild und Erzeugung vollständig bekannt und nachvollziehbar sind und der damit die kontrollierte Vergleichsgrundlage für die weitere Untersuchung bildet. Zugleich ist die dem Vorgehen innewohnende Einschränkung zu benennen: Ein synthetischer Datensatz kann die Vielfalt und die unvorhergesehenen Unregelmäßigkeiten realer Datenbestände nur annähern, weshalb die aus ihm gewonnenen Ergebnisse unter diesem Vorbehalt zu interpretieren und im Rahmen der Limitationen in Kapitel~\ref{ch:fazit} erneut aufzugreifen sind \autocite{liSyntheticDataGeneration2023,goyalSystematicReviewSynthetic2024}. Auf der so geschaffenen Grundlage kann im folgenden Kapitel~\ref{ch:implementierung} die technische Umsetzung der Verarbeitungsumgebung beschrieben werden, welche die hier erzeugten Daten sowohl der klassischen Pipeline als auch den \ac{LLM}-gestützten Verfahren zuführt.
